\section{Introduction}\label{sec:intro}

A benchmark number is only as informative as the split that produced it. When a dataset is
released as pre-computed spectrogram images, one image corresponds to a short window of one
recording, and a model trained on images from a recording has already seen the acoustics of
that recording: instrument timbre, room, performer, and microphone. A split over images can
therefore test within-recording recognition rather than generalisation to unseen recordings.
This is a case of non-independence between training and test, one category in the leakage
taxonomy of Kapoor and Narayanan~\cite{kapoor2023leakage}. The magnitude is an empirical
question; our paired audit measures it without changing the backbone.

Music information retrieval, or MIR, has dealt with this problem for a long time. Splitting by recording is
used in several audio benchmarks, for example urban sound, instrument and genre corpora, and the
split unit is usually stated explicitly; Salamon et al.~\cite{salamon2014urbansound} fixed a recording-level split for
UrbanSound8K in 2014, keeping slices of one recording within one fold so that evaluation is
recording-disjoint and comparable across papers. Yet datasets
that distribute features instead of audio make the split unit easy to lose, and multi-library
benchmarks, which pool heterogeneous sources under one evaluation protocol, make the loss hard
to notice: each library carries its own release format, its own file naming, and its own notion
of what one ``item'' is.

This paper audits one such framework. CCMusic \citep{zhou2025ccmusic} is a recent Chinese-music
resource whose associated article documents six datasets and evaluates all six under one
framework; the public release carries thirteen libraries under one publishing organisation, of
which those six are a subset. It distributes chroma, CQT and MEL features as pre-computed
images, and its evaluation framework states a split ratio for each collection --- 8:1:1 for
CTIS after 2-second segmentation, a recording-level 79:10:10 for Guzheng Tech99, 6:2:2 for
ErhuPT --- naming the split unit only in the recording-level case. We examine a
slice-level interpretation supported by the released format: MEL images are named by recording
and window index, and a naive 8:1:1 split over image filenames places windows of the same
recording on both sides. The released archives make this reconstruction possible, but do not
establish the split actually used by the original authors.

Our contribution is a measurement, not a discovery of a new phenomenon. The framework's own
documentation leaves three quantities unstated, and we estimate each of them. The first is the
size of the A1$-$A0 contrast under equal-volume replacement with matched training-pool
recording counts, while acknowledging that A1 introduces additional pool recordings. The second is whether a recording-level split is feasible for these libraries at their
stated class counts, given that splitting by recording consumes recordings that would otherwise
train the model. The third is whether the pipeline yields small contrasts in near-single-slice controls,
where homologous exposure is absent or very limited by construction.

To answer them we rebuilt the five libraries that release MEL slices --- 110{,}415 released
images from 7{,}827 recordings across 252 classes, of which 109{,}831 images from 7{,}756
recordings across 234 classes survive a minimum-recording rule --- under a recording-level
three-way split, and trained two standard backbones, AlexNet and Swin-T, under three main arms.
A CTIS-only fourth arm, C1, was added to check equal-volume non-test injection. The control arm, A0, samples a fixed number of
slices per class from the training recordings and never touches a test recording. The exposure
arm, A1, injects the entire held-out Q pool by equal replacement. Many pool slices come from
recordings also represented in the test set, but not all do; realised overlap is reported in
Supplementary Material S1. The number of training slices is held fixed and the
number of training recordings drawn from the training pool is matched. A1 differs from A0 in
whose content reaches training and in the additional pool recordings, whose counts are reported
per class; the contrast is therefore consistent with an advantage associated with
within-recording exposure, but it does not isolate one. A2 provides an illustrative self-test
score under our reconstruction of a slice-level split over the pooled corpus; it is neither a
reproduction of the original authors' unknown split files nor the original benchmark score.

Three results organise the paper. Conditional on the fixed split and the trained seeds, the
exposure arm gains $+2.02$ points of macro-F1 in the primary library, with a 95\% CI of $+0.83$ to $+3.35$, and $+1.70$ under
Swin-T with CI $+1.28$ to $+2.84$. Both intervals exclude zero. The Swin-T interval
lies wholly below our three-point smallest effect of interest;
the AlexNet interval includes three points, so that magnitude is not excluded. A CTIS-only
control arm that supplies the same volume of training slices from non-test content, C1, moves
the point estimate by $+0.05$ and $+0.62$ points only; measured against C1, the exposure arm
still gains $+1.96$ points $[+0.74,+2.99]$ under AlexNet and $+1.07$ $[+0.23,+1.96]$ under
Swin-T, so the difference is smaller than the A0 contrast, though its paired interval excludes zero under
the fixed split and the trained seeds. In four other libraries
the same manipulation moves macro-F1 by $+8.4$ to $+28.7$ points. In the three negative-control
libraries --- near-single-slice collections --- five of six contrasts include zero. The sixth,
chest falsetto under Swin-T, is $-2.53$ $[-5.00,-0.21]$; its three seed-specific deltas are
$+5.57$, $-2.62$ and $-10.54$ points, so the sign is not stable across seeds. With six
contrasts and no correction for multiplicity, these controls check the pipeline, not the
mechanism in isolation, and they do not establish equivalence.

We are careful about what these numbers license. They do not show that CCMusic is
irreproducible, nor that its authors made an error that invalidates published results: the
audited article is a database paper, and its evaluation framework is one section of it. They do
show a small positive equal-replacement contrast in CTIS and larger contrasts in four additional
libraries, conditional on our fixed reconstruction and training recipe. They also give the
arithmetic needed to judge whether a recording-disjoint split is feasible at the stated
per-class requirements, rather than treating feature images as independent examples.

The remainder of the paper is organised as follows. Section~\ref{sec:related} positions the
audit against leakage taxonomy, audio-dataset practice, and dataset audits in adjacent fields.
Section~\ref{sec:framework} states what the audited framework says and what its released files
imply. Section~\ref{sec:materials} describes the five libraries, the manifest, and the
capacity calculation. Section~\ref{sec:protocol} defines the arms, the sealing protocol, and
the statistical plan. Section~\ref{sec:results} reports the main effect, the random-split arm,
the negative controls, and a cross-instrument experiment on one library.
Section~\ref{sec:discussion} discusses what the size of the effect depends on and what a
reporting standard should require. Section~\ref{sec:conclusion} concludes.
